\section{Related Work}
\label{sec:related-work}

\paragraph{Coverage-guided fuzzing}
American Fuzzy Lop~\cite{zalewski2014afl} and its successor AFL++~\cite{fioraldi2020aflplusplus}
popularized mutation-based, coverage-guided fuzzing: a target is
instrumented to report which code ran, and the fuzzer retains and mutates
inputs that reach new coverage. Man\`es et al.'s survey~\cite{manes2021fuzzing}
gives a taxonomy of this and related approaches. Coverage-guided fuzzers are
highly effective when instrumentation is available, but for a structured
input format their byte-level mutations mostly produce syntactically invalid
inputs that are rejected before any interesting code runs, and none of this
family applies when the target cannot be recompiled with instrumentation --
the setting this paper addresses.

\paragraph{Grammar-based and generation-based fuzzing}
Generating inputs directly from a grammar, rather than mutating bytes, is a
much older idea: Purdom's sentence generator~\cite{purdom1972sentence}
produces short strings that exercise every grammar production, originally
for testing and debugging parsers themselves. More recent tools combine a
grammar with coverage feedback: Superion~\cite{wang2019superion} extends a
greybox fuzzer with grammar-aware mutation operators so mutated inputs
remain syntactically valid, and Nautilus~\cite{aschermann2019nautilus}
similarly pairs a grammar-based generator with coverage-guided mutation to
reach deeper parser states. Both still assume coverage instrumentation is
available to decide which grammar-derived inputs to keep; our setting
removes that assumption entirely; the only feedback available is the kind a
black-box harness can itself produce, so the loop steers the grammar-based
\emph{generator}, not a corpus of retained inputs.

\paragraph{Property-based testing}
QuickCheck~\cite{claessen2000quickcheck} introduced property-based testing:
a developer writes generators for values of a given shape and a property
that should hold for all of them, and the tool searches for a
counterexample. Hypothesis~\cite{maciver2019hypothesis}, which this work
builds directly on, is a Python implementation of the same idea, with
composable strategies that can mirror a formal grammar. In this line of work
the generator is authored once by a human and is not revised in response to
how the target under test actually behaves. Vikram et
al.~\cite{vikram2023pbt} show that LLMs can write property-based tests,
including their input generators, from API documentation, and separate a
test's validity from its usefulness; we ask the corresponding question for a
grammar-based generator and a black-box target.

\paragraph{LLMs for fuzzing and test generation}
LLMs have recently been used to generate or mutate fuzzing inputs directly.
TitanFuzz~\cite{deng2023titanfuzz} uses an LLM to generate and mutate
programs for fuzzing deep-learning library APIs, relying on the LLM's prior
exposure to real usage examples of those APIs rather than a formal grammar.
Fuzz4All~\cite{xia2024fuzz4all} generalizes this to arbitrary
languages/systems by using an LLM both to synthesize a distilled input
prompt from documentation and examples and to generate and mutate inputs
from it. Ackerman and Cybenko~\cite{ackerman2023fuzzingparsers} have an LLM
read a natural-language format specification and generate parser inputs as
seeds for a mutation fuzzer. CodaMOSA~\cite{lemieux2023codamosa} takes a related but distinct
approach for unit-test generation: search-based test generation runs until
coverage plateaus, at which point an LLM is asked for example test cases to
help the search escape the plateau -- i.e., the LLM is invoked \emph{because}
coverage feedback is available and has stalled. Closest to our loop,
G2Fuzz~\cite{zhang2025g2fuzz} has an LLM synthesize and mutate input
generators for non-textual formats and pairs them with a coverage-guided
mutation fuzzer, and ELFuzz~\cite{chen2025elfuzz} evolves generation-based
fuzzers with an LLM under coverage guidance. Both revise generators, as we
do, but select them by coverage. Our setting differs from all of these
along the same axis: we assume no coverage signal is available at any
point, only the parser's own accept/reject decision and sanitizer output,
and the LLM's job is not to produce one-off inputs or test cases but to
revise a reusable, executable \emph{generator} against a formal grammar it is
also given directly, iterating on the same three coverage-free proxies each
round.

\paragraph{Testing JSON parsers}
JSONTestSuite~\cite{seriot2016jsontestsuite} labels test documents as ones a
parser must accept, must reject, or may handle either way under RFC
8259~\cite{rfc8259}, and Möller et al.~\cite{moller2024crosslanguage}
differentially test JSON parsers across languages and report many semantic
discrepancies. Both start from known inputs. We study a generator that must
reach such inputs on its own, and Section~\ref{sec:evaluation-parson} shows
that the rejections on our second target fall into these suites' categories.

\paragraph{Positioning} This paper sits at the intersection of the second
and fourth lines above: like Superion and Nautilus, it uses a grammar to
keep generated inputs structurally valid; unlike them, it has no coverage
instrumentation to fall back on, and unlike TitanFuzz, Fuzz4All, and
CodaMOSA, the LLM never sees coverage data, execution traces, or usage
examples, and unlike G2Fuzz and ELFuzz, no coverage signal decides which
generator survives -- only the grammar text and the same three black-box proxies
(acceptance rate, structural diversity, rejection signatures) a human
maintaining a Hypothesis strategy by hand would have to reason about
manually.
